\section{Introduction}

\label{sec:introduction}
% ══════════════════════════════════════════════════════════════════════════════

Biodiversity monitoring at global scales requires automated species
identification systems capable of operating across the full breadth of
taxonomic diversity \citep{christin2019applications}. Camera traps,
acoustic sensors, citizen science platforms, and environmental DNA (eDNA)
sampling generate petabytes of observational data annually, far exceeding
manual identification capacity \citep{steenweg2017scaling}.

Deep learning has revolutionized image-based species identification, with
convolutional neural networks achieving expert-level performance on
well-studied taxa such as birds \citep{van2015building}, plants
\citep{goeau2022overview}, and mammals \citep{norouzzadeh2018automatically}.
However, these successes depend on large labeled training datasets---a
critical limitation given that:

\begin{enumerate}[label=(\roman*)]
  \item The majority of species have fewer than 50 available images in
    public databases \citep{lorieul2022overview}.
  \item Newly discovered species (approximately 18{,}000 per year
    \citep{costello2013can}) have no training data at all.
  \item Rare and cryptic species---often those of greatest conservation
    concern---are systematically underrepresented in image databases.
  \item Intraspecific visual variation (sexual dimorphism, juvenile
    plumage, seasonal morphs) multiplies the effective number of
    required examples.
\end{enumerate}

Few-shot learning addresses this data scarcity by training models to
recognize novel classes from very few examples
\citep{wang2020generalizing}. While few-shot methods have advanced
rapidly in computer vision, their application to ecological species
identification introduces unique challenges:

\begin{itemize}
  \item \textbf{Fine-grained discrimination}: Species within the same
    genus can differ by subtle visual features (wing venation, petal
    shape, scale patterns) that require high-resolution feature
    extraction.
  \item \textbf{High inter-class similarity}: Closely related species
    form dense clusters in feature space, making prototype-based
    methods fragile.
  \item \textbf{Domain shift}: Training images from curated databases
    differ substantially from field photos taken by citizen scientists
    under variable conditions.
  \item \textbf{Hierarchical structure}: Biological taxonomy provides a
    natural hierarchy that standard few-shot methods ignore.
\end{itemize}

\subsection{Contributions}

We introduce \textbf{TaxoNet}, a few-shot species recognition system
with the following contributions:

\begin{enumerate}
  \item A \textbf{taxonomy-aware prototypical network} that constructs
    multi-rank prototypes (family, genus, species) and fuses them
    through learned attention, enabling few-shot identification to
    leverage phylogenetic structure.

  \item A \textbf{hierarchical meta-learning} training strategy with
    taxonomic episode sampling that ensures the model learns
    discriminative features at every rank.

  \item A \textbf{domain-adaptive pre-training} pipeline using 28
    million ecological images from iNaturalist, GBIF, and museum
    collections.

  \item Integration with \textbf{geographic and phenological priors}
    that improve practical field identification accuracy by
    constraining the candidate species set.

  \item Comprehensive evaluation across \textbf{four major taxonomic
    groups} (birds, insects, plants, marine organisms) with separate
    validation on \textbf{citizen science field data}.
\end{enumerate}

\subsection{Antje Worring, Zach Kelling Context}

This work supports Antje Worring, Zach Kelling's mission of open, decentralized
science infrastructure. TaxoNet is designed for deployment on the Zoo
citizen science platform and integrates with the Foundation's
decentralized data pipelines. Model weights and training code are
released under the Apache 2.0 license, and the evaluation benchmark
is maintained as a community resource through the Zoo Open Science
Network at \texttt{zips.zoo.ngo}.


% ══════════════════════════════════════════════════════════════════════════════
